%% kapa: exam=1 round=33 year=2022 period=1 form=A subject=real_estate subject_ko=부동산학원론 seq=14 number=94
%% source: Q-Net 감정평가사 기출문제 (cst00309) artlSeq=5209997 file=raw/1008/5209997/2022년도 제33회 감정평가사 1차시험 1교시(민법,경제학,부동산학).pdf
%% conversion: kordoc 4.24.51 -> tools/gen1.py (auto); review=proofread-sonnet; flags=none
\begin{question}{94}
부동산투자회사법상 부동산투자회사에 관한 설명으로 옳은 것은?
\begin{choices}
  \item 최저자본금준비기간이 지난 위탁관리 부동산투자회사의 자본금은 70억원 이상이 되어야 한다.
  \item 자기관리 부동산투자회사의 설립자본금은 3억원 이상으로 한다.
  \item 자기관리 부동산투자회사에 자산운용 전문인력으로 상근하는 감정평가사는 해당분야에 3년 이상 종사한 사람이어야 한다.
  \item 최저자본금준비기간이 끝난 후에는 매 분기 말 현재 총자산의 100분의 80 이상이 부동산(건축 중인 건축물 포함)이어야 한다.
  \item 위탁관리 부동산투자회사는 해당 연도 이익을 초과하여 배당할 수 있다.
\end{choices}
\end{question}
\answer{5}
